% In IRMF.tex, insert % In IRMF.tex, insert % In IRMF.tex, insert % In IRMF.tex, insert \input{irmf_component_proofs} immediately after \input{IMF}.
\section{Finite-sample bounds for the IRMF components}\label{irmfproof:section}

At the first fit, the estimation error is the empirical score divided by
the population curvature, up to a smaller remainder. At the second fit,
part of that first error has already been subtracted from the data.
We show how to subtract its contribution from the second score and repeat
the argument at each scale. For a window containing \(m\) equally weighted
observations, the leading error at that scale is of order \(\sqrt{z/m}\)
and, under the conditions below, the remainder is of order \(z/m\), where \(z=\xx+\log(10n k_{\max})\)
at confidence \(1-\ex^{-\xx}\).

Throughout, the \(n\) grid points, the signal and the bandwidths are fixed,
\(\Data_u=\Signal_u+\eps_u\) with i.i.d.\ errors of finite first moment,
and the kernel weights are nonnegative with positive row sums. All maxima
over \(t\) are on this grid.

\subsection{The first component}\label{irmfproof:first}

In \eqref{miuygtd5r6wguwhd7r54we} let the contrast be convex, coercive and
\(C^3\), with \(|\rhohu'|\le L_1\), \(0\le\rhohu''\le L_2\), and
\(|\rhohu'''|\le L_3\). Write \(h=h_1\), \(Q_t^2=\sum_u\Kern_h(t-u)^2\),
\(\kappa_t=\max_u\Kern_h(t-u)\), \(\DFN_{\min}=\min_t\sqrt{\IF_t}>0\),
and \(z_1=\xx+\log(4n)\) for \(\xx>0\). The score, curvature and third
derivative in Lemma~\ref{LRMFloc} have respective bounds
\begin{align*}
 \CONSTi_1&=\max_t L_1(Q_t\sqrt{2z_1}+2\kappa_t z_1/3)/\DFN_t,\\
 \CONSTi_2&=\max_t L_2(Q_t\sqrt{2z_1}+\kappa_t z_1/3)/\DFN_t,
 &\CONSTi_3&=\max_t L_3N_h(t)/\DFN_t^2 .
\end{align*}

\begin{proposition}\label{irmfproof:first-bound}
If \(\delta_*=\CONSTi_2/\DFN_{\min}<1\) and
\(\CONSTi_1\CONSTi_3/[\DFN_{\min}(1-\delta_*)^2]<4/9\), then with
probability at least \(1-\ex^{-\xx}\),
\begin{equation}\label{irmfproof:first-expansion}
 \max_t\DFN_t|\IMFt_t-\IMFs_t-\IF_t^{-1}\nabla_t|
 \le \frac{1}{\DFN_{\min}}\left\{
 \frac{\CONSTi_1\CONSTi_2}{1-\delta_*}
 +\frac{9\CONSTi_1^2\CONSTi_3}{8(1-\delta_*)^3}\right\}.
\end{equation}
\end{proposition}

\begin{proof}
Fix \(t\), set \(\xs=\IMFs_t\), and use the empirical contrast \(\fp\)
in \eqref{fjiuy9iw3ojlggeefr}. The population score equation gives
\(\fp'(\xs)=-\nabla_t\). The centered score summands have total variance
at most \(L_1^2 Q_t^2\) and maximum size \(2L_1\kappa_t\); the corresponding
curvature bounds are \(L_2^2 Q_t^2\) and \(L_2\kappa_t\).
Bernstein's inequality gives \(|\nabla_t|\le\DFN_t\CONSTi_1\) and
\(|\IFP-\IF_t|\le\DFN_t\CONSTi_2\) for all \(t\), with failure probability
at most \(4n\ex^{-z_1}\), where \(\IFP=\fp''(\xs)\).
The third-derivative bound gives
\(\max(\sup|\fs'''|,\sup|\fp'''|)\le\DFN_t^2\CONSTi_3\), so
\(|\fs''(y)-\fs''(\xs)|\le\DFN_t^2\CONSTi_3|y-\xs|\).

In curvature units \(v=\DFN_t(y-\xs)\), put \(g(v)=\fp(y)\),
\(c=g''(0)\ge1-\delta_*\), \(b=|g'(0)|/c\), and
\(\tau=\CONSTi_3/\DFN_t\). Taylor's inequality gives
\(|g'(v)-g'(0)-cv|\le\tau v^2/2\). Since \(\tau b/c<4/9\), the
derivative has opposite signs at \(\pm3b/2\) if \(b>0\), and curvature
stays positive between them. Thus its zero \(v_\circ\) satisfies
\(|v_\circ|\le3b/2\), also when \(b=0\). At that zero,
\(|v_\circ+g'(0)/c|\le9\tau b^2/(8c)\). This bounds the error using
\(\IFP^{-1}\nabla_t\). Replacing \(\IFP\) by \(\IF_t\) costs at most
\(\CONSTi_1\CONSTi_2/[\DFN_t(1-\delta_*)]\) in curvature units.
Use \(b\le\CONSTi_1/(1-\delta_*)\) and maximize over \(t\). The same localization gives
\(\max_t\DFN_t|\IMFt_t-\IMFs_t|\le3\CONSTi_1/[2(1-\delta_*)]\).
\end{proof}

\subsection{How the first error enters the later fits}\label{irmfproof:sec:recursive}

Take the Gaussian-smoothed absolute loss \(\rhohu=\rhohu_H\) of
Lemma~\ref{Lrhohu}. Write \(\KS_{tu}^{(k)}=\Kern_{h_k}(t-u)/N_k(t)\)
for the weights in \eqref{dfikkj3ujvyhwedrgg} and
\(\alpha_k=\max_{t,u}\KS_{tu}^{(k)}\).
The true component \(\IMFs^{(k)}\) uses the law of the actual residual
\(\Data^{(k)}\). For the proof we also need a deterministic reference:
start at \(\Signal_{\rhohum}^{(1)}=\Signal\), fit each row using
\(\rhohum(y)=\E\rhohu(y+\eps_1)\), then subtract that fit,
\[
 \IMFs_{\rhohum,t}^{(k)}
  =\arginf_s\sum_u\KS_{tu}^{(k)}\rhohum(\Signal_{\rhohum,u}^{(k)}-s),
 \qquad \Signal_{\rhohum}^{(k+1)}=\Signal_{\rhohum}^{(k)}-\IMFs_{\rhohum}^{(k)}.
\]
The two targets coincide at stage one. On the reference path, define
\(y_{ktu}=\Signal_{\rhohum,u}^{(k)}+\eps_u-\IMFs_{\rhohum,t}^{(k)}\),
\(g_{ktu}=\rhohu''(y_{ktu})\), and \(f_{kt}=\sum_u\KS_{tu}^{(k)}\E g_{ktu}\).
The score term \(b_k\) and the matrix \(V_k\) that transmits a change
in the input are
\[
 b_{kt}=\frac{\sum_u\KS_{tu}^{(k)}\rhohu'(y_{ktu})}{f_{kt}},
 \qquad V_k[t,u]=\frac{\KS_{tu}^{(k)}\E g_{ktu}}{f_{kt}} .
\]
In particular \(b_{1t}=\IF_t^{-1}\nabla_t\). At the next fit the leading
error is \(b_2-V_2b_1\). In general, accumulate the preceding leading
errors in \(\Lambda_k\) and set
\begin{equation}\label{irmfproof:eq:lambda}
 \Lambda_1=0,\qquad \mathcal L_k=b_k-V_k\Lambda_k,\qquad
 \Lambda_{k+1}=\Lambda_k+\mathcal L_k .
\end{equation}

\paragraph{Conditions.}
Assume \(f_{kt}\ge\lambda>0\) and
\(\alpha_j\le C_\alpha\alpha_k\hband^{-(k-j)}\) for \(j\le k\),
with \(\hband>1\), \(C_\alpha\ge1\). Write
\(B_{m:l}=(\Id-V_m)\cdots(\Id-V_l)\), with empty product \(\Id\).
We need stability of every such product and limited variation of curvature:
\begin{equation}\label{irmfproof:eq:dispersion}
 \max_{l\le m}\|B_{m:l}\|_\infty\le C_B,\qquad
 d=\max_{k,t}\frac{\sum_u\KS_{tu}^{(k)}\E|g_{ktu}-\E g_{ktu}|}{f_{kt}}
       <\frac{\hband-1}{C_B C_\alpha},\qquad C_B\ge1 .
\end{equation}
Here \(\|\cdot\|_\infty\) is the maximum absolute row sum for matrices and
the maximum absolute entry for vectors.

\begin{theorem}\label{irmfproof:thm:main}
Put \(u_k=\sqrt{z\alpha_k}\) and \(\alpha_{\max}=\max_k\alpha_k\).
There are constants \(C,\eta>0\) depending only on \(H,\lambda,\hband,
C_\alpha,C_B,d\) such that, if \(z\alpha_{\max}\le\eta\), then with
probability at least \(1-\ex^{-\xx}\), simultaneously for \(1\le k\le k_{\max}\),
\begin{equation}\label{irmfproof:eq:expansion}
 \|\mathcal L_k\|_\infty\le C u_k,\qquad
 \|\IMFt^{(k)}-\IMFs_{\rhohum}^{(k)}-\mathcal L_k\|_\infty\le C u_k^2 .
\end{equation}
\end{theorem}

\begin{proof}
Write \(e_k=\IMFt^{(k)}-\IMFs_{\rhohum}^{(k)}\) and
\(E_k=\sum_{j<k}e_j\). The actual input is exactly
\(\Data^{(k)}=\Signal_{\rhohum}^{(k)}+\eps-E_k\).
We use \(C\) for constants with the dependence stated above and put \(u_0=0\).

\emph{Expand the score once.}
The fitted location solves
\(\sum_u\KS_{tu}^{(k)}\rhohu'(y_{ktu}-E_k(u)-e_k(t))=0\).
Define \(Z_k[t,u]=\KS_{tu}^{(k)}(g_{ktu}-\E g_{ktu})/f_{kt}\).
Taylor expansion at \(y_{ktu}\), followed by division by \(f_{kt}\), gives
\begin{equation}\label{irmfproof:eq:recursion}
 e_k=b_k-V_kE_k+w_k,\qquad
 \|w_k+Z_kE_k\|_\infty
 \le \|Z_k\mathbf{1}\|_\infty\|e_k\|_\infty+
 \frac{L_3}{2\lambda}(\|E_k\|_\infty+\|e_k\|_\infty)^2 .
\end{equation}
Here \(Z_k\mathbf{1}\) is the row sum of the centered curvature. It is small
even when the sum of the absolute entries in that row is appreciable.

\emph{Control the terms that share observations.}
Since \(\sum_u(\KS_{tu}^{(k)})^2\le\alpha_k\), Bernstein's inequality
applied to the scores, curvatures and absolute centered curvatures gives,
simultaneously, \(\|b_k\|_\infty+\|Z_k\mathbf{1}\|_\infty\le C u_k\)
and \(\|Z_k\|_\infty\le d+C u_k\) when \(u_k\le1\).
For the input \(\Signal_{\rhohum}^{(k)}+\eps\), let \(s\) be the absolute
score at the reference fit. Here \(s\le C u_k\) and the empirical curvature
is at least \(3\lambda/4\) for small \(\eta\). Since it is \(L_3\)-Lipschitz,
it stays above \(\lambda/2\) within \(2s/\lambda\), which brackets the fit.
A local fit is monotone in each
observation and translation equivariant, hence 1-Lipschitz in the sup norm.
Consequently \(\|e_k\|_\infty\le C u_k+\|E_k\|_\infty\).

Iterating \eqref{irmfproof:eq:lambda} gives
\(\Lambda_k=\sum_{j<k}B_{k-1:j+1}b_j\). The geometric schedule then gives
\(\|\Lambda_k\|_\infty\le C_L u_{k-1}\).
The same sum grouped by observation gives
\(\Lambda_k(u)=\sum_v\gamma_{uv}(\eps_v)\), with each summand centered.
Indeed \(\gamma_{uv}\) is a sum of centered scores at \(v\), bounded in absolute value by
\(G_{uv}=2\lambda^{-1}\sum_{j<k}\sum_s|B_{k-1:j+1}[u,s]|\KS_{sv}^{(j)}\).
The product bound and the weight schedule imply
\(\max_{u,v}G_{uv}\le G_\infty\alpha_{k-1}\) and
\(\max_u\|G_{u,\cdot}\|_2\le C\sqrt{\alpha_{k-1}}\).
In \(Z_k\Lambda_k\), the terms \(u=v\) thus contribute at most
\(C\alpha_{k-1}\). The terms \(u\ne v\) are centered in each argument;
the sum of their squared envelopes is at most \(C\alpha_k\alpha_{k-1}\).
The quadratic-sum bound in the companion proof gives
\(\|Z_k\Lambda_k\|_\infty\le C u_k^2\).%
\footnote{For independent coordinates and separately centered kernels
\(|h_{uv}|\le H_{uv}\), the off-diagonal sum is at most
\(32\ex z(\sum H_{uv}^2)^{1/2}\) except with probability \(\ex^{-z}\)
for \(z\ge2\). A proof by symmetrization and explicit constants for the
component bounds are in Appendices A and B of the
\href{https://github.com/mkuziuk/imf/blob/ade89c8fd39ff9b3461ae68b33227ae08f9865a4/overleaf/irmf_all_components_bound.pdf}{companion proof}.}
The first three Bernstein bounds and this last bound cost at most
\(6n k_{\max}\ex^{-z}\le\ex^{-\xx}\) in total.

\emph{Close the induction.}
Set \(\Theta_k=E_k-\Lambda_k\). Choose
\(d<\delta<(\hband-1)/(C_B C_\alpha)\) and reduce \(\eta\) until
\(\|Z_k\|_\infty\le\delta\).
If \(\|\Theta_j\|_\infty\le C_\Theta u_{j-1}^2\) through stage \(k\), then
\(\|E_j\|_\infty\le(C_L+1)u_{j-1}\) provided \(C_\Theta\sqrt\eta\le1\).
Splitting \(Z_j E_j=Z_j\Lambda_j+Z_j\Theta_j\) in
\eqref{irmfproof:eq:recursion} now gives
\(\|w_j\|_\infty\le C_0 u_j^2+\delta C_\Theta u_{j-1}^2\), with \(C_0\)
independent of \(C_\Theta\). Since
\(\Theta_{k+1}=\sum_{j\le k}B_{k:j+1}w_j\), summing the two geometric
series gives
\[
 \|\Theta_{k+1}\|_\infty
 \le\left\{\frac{C_B C_\alpha C_0}{1-\hband^{-1}}
       +\frac{C_B C_\alpha\delta}{\hband-1}C_\Theta\right\}u_k^2 .
\]
The coefficient of \(C_\Theta\) is less than one. Choose \(C_\Theta\) large enough to
close this bound, then choose \(\eta\) as above. The induction starts at
\(E_1=\Lambda_1=0\). Finally \(e_k-\mathcal L_k=\Theta_{k+1}-\Theta_k\)
and \(u_{k-1}^2\le(C_\alpha/\hband)u_k^2\), which proves both bounds.
\end{proof}

\subsection{Returning to the true IMF}\label{irmfproof:sec:targets}

The first-stage target is already \(\IMFs_t\). At later stages the shift depends on how earlier fits changed the residuals.
In the expected score, contributions from distinct observations cancel.
Write \(z_0=\log(10n k_{\max})\), \(\xx_*=\eta/\alpha_{\max}-z_0>0\), and
assume \(\mu_4=\E|\eps_1|^4<\infty\). Let \(m_{\rhohu}\) solve
\(\E\rhohu'(\eps_1-m_{\rhohu})=0\) and put
\(M=2^{k_{\max}+1}(\|\Signal\|_\infty+(n\mu_4)^{1/4}+|m_{\rhohu}|)\).
With the induction constants above, set \(A=C_L+1\) and, for \(k\ge2\),
\begin{align*}
 U_k&=C_\Theta(z_0+1)\alpha_{k-1}+(M+C_B k_{\max}/\lambda)\ex^{-\xx_*/2},\\
 T_k&=A^2(z_0+1)\alpha_{k-1}+M^2\ex^{-\xx_*/2},\\
 \beta_k&=L_2(G_\infty\alpha_{k-1}+U_k)+L_3T_k/2,\qquad \beta_1=0 .
\end{align*}
\begin{corollary}\label{irmfproof:cor:population}
If \(L_3\sqrt{T_k}\le\lambda/4\) and \(\beta_k\le\lambda^2/(8L_3)\)
at each \(k\ge2\), the event of Theorem~\ref{irmfproof:thm:main} gives
\[
 \|\IMFt^{(k)}-\IMFs^{(k)}\|_\infty\le C(u_k+u_k^2)+2\beta_k/\lambda .
\]
If both exponential tail terms in \(U_k,T_k\) are at most \(\min_j\alpha_j\),
the shift is \(O(z\alpha_k)\). Thus for \(\alpha_k\asymp(nh_k)^{-1}\) the
simultaneous rate is \(\sqrt{z/(nh_k)}+z/(nh_k)\).
\end{corollary}

\begin{proof}
The induction gives \(\E\|\Theta_k\|_\infty\le U_k\) and
\(\E\|E_k\|_\infty^2\le T_k\): integrate its bounds over
\(0<\xx\le\xx_*\), using \(\int_0^\infty\ex^{-s}ds=1\).
Above the cutoff use Cauchy--Schwarz on an event of probability at most
\(\ex^{-\xx_*}\). Indeed an empirical fit lies between its extreme inputs
and a reference fit between those extremes shifted by \(m_{\rhohu}\);
iterating these brackets gives \(\|\|E_k\|_\infty\|_{L^4}\le M\).
Also \(\|\Lambda_k\|_\infty\le C_B k_{\max}/\lambda\) pathwise.

At \(\IMFs_{\rhohum,t}^{(k)}\) the true population score is
\(\sum_u\KS_{tu}^{(k)}\E\rhohu'(y_{ktu}-E_k(u))\).
Expand once about \(y_{ktu}\). The constant term sums to zero. In the linear term \(\E[g_{ktu}\Lambda_k(u)]\), independence and centering
cancel every \(\gamma_{uv}(\eps_v)\) except \(v=u\), leaving at most
\(L_2 G_\infty\alpha_{k-1}\). The remaining linear and quadratic terms
are bounded by \(L_2 U_k\) and \(L_3T_k/2\). Thus the score has magnitude
at most \(\beta_k\). The population loss has curvature at least
\(\lambda-L_3\sqrt{T_k}\) at the reference fit and at least \(\lambda/2\) within
\(2\beta_k/\lambda\). The score therefore has its zero within \(2\beta_k/\lambda\) of the reference fit.
Adding this displacement to \eqref{irmfproof:eq:expansion} proves the claim.
\end{proof}
 immediately after \input{IMF}.
\section{Finite-sample bounds for the IRMF components}\label{irmfproof:section}

At the first fit, the estimation error is the empirical score divided by
the population curvature, up to a smaller remainder. At the second fit,
part of that first error has already been subtracted from the data.
We show how to subtract its contribution from the second score and repeat
the argument at each scale. For a window containing \(m\) equally weighted
observations, the leading error at that scale is of order \(\sqrt{z/m}\)
and, under the conditions below, the remainder is of order \(z/m\), where \(z=\xx+\log(10n k_{\max})\)
at confidence \(1-\ex^{-\xx}\).

Throughout, the \(n\) grid points, the signal and the bandwidths are fixed,
\(\Data_u=\Signal_u+\eps_u\) with i.i.d.\ errors of finite first moment,
and the kernel weights are nonnegative with positive row sums. All maxima
over \(t\) are on this grid.

\subsection{The first component}\label{irmfproof:first}

In \eqref{miuygtd5r6wguwhd7r54we} let the contrast be convex, coercive and
\(C^3\), with \(|\rhohu'|\le L_1\), \(0\le\rhohu''\le L_2\), and
\(|\rhohu'''|\le L_3\). Write \(h=h_1\), \(Q_t^2=\sum_u\Kern_h(t-u)^2\),
\(\kappa_t=\max_u\Kern_h(t-u)\), \(\DFN_{\min}=\min_t\sqrt{\IF_t}>0\),
and \(z_1=\xx+\log(4n)\) for \(\xx>0\). The score, curvature and third
derivative in Lemma~\ref{LRMFloc} have respective bounds
\begin{align*}
 \CONSTi_1&=\max_t L_1(Q_t\sqrt{2z_1}+2\kappa_t z_1/3)/\DFN_t,\\
 \CONSTi_2&=\max_t L_2(Q_t\sqrt{2z_1}+\kappa_t z_1/3)/\DFN_t,
 &\CONSTi_3&=\max_t L_3N_h(t)/\DFN_t^2 .
\end{align*}

\begin{proposition}\label{irmfproof:first-bound}
If \(\delta_*=\CONSTi_2/\DFN_{\min}<1\) and
\(\CONSTi_1\CONSTi_3/[\DFN_{\min}(1-\delta_*)^2]<4/9\), then with
probability at least \(1-\ex^{-\xx}\),
\begin{equation}\label{irmfproof:first-expansion}
 \max_t\DFN_t|\IMFt_t-\IMFs_t-\IF_t^{-1}\nabla_t|
 \le \frac{1}{\DFN_{\min}}\left\{
 \frac{\CONSTi_1\CONSTi_2}{1-\delta_*}
 +\frac{9\CONSTi_1^2\CONSTi_3}{8(1-\delta_*)^3}\right\}.
\end{equation}
\end{proposition}

\begin{proof}
Fix \(t\), set \(\xs=\IMFs_t\), and use the empirical contrast \(\fp\)
in \eqref{fjiuy9iw3ojlggeefr}. The population score equation gives
\(\fp'(\xs)=-\nabla_t\). The centered score summands have total variance
at most \(L_1^2 Q_t^2\) and maximum size \(2L_1\kappa_t\); the corresponding
curvature bounds are \(L_2^2 Q_t^2\) and \(L_2\kappa_t\).
Bernstein's inequality gives \(|\nabla_t|\le\DFN_t\CONSTi_1\) and
\(|\IFP-\IF_t|\le\DFN_t\CONSTi_2\) for all \(t\), with failure probability
at most \(4n\ex^{-z_1}\), where \(\IFP=\fp''(\xs)\).
The third-derivative bound gives
\(\max(\sup|\fs'''|,\sup|\fp'''|)\le\DFN_t^2\CONSTi_3\), so
\(|\fs''(y)-\fs''(\xs)|\le\DFN_t^2\CONSTi_3|y-\xs|\).

In curvature units \(v=\DFN_t(y-\xs)\), put \(g(v)=\fp(y)\),
\(c=g''(0)\ge1-\delta_*\), \(b=|g'(0)|/c\), and
\(\tau=\CONSTi_3/\DFN_t\). Taylor's inequality gives
\(|g'(v)-g'(0)-cv|\le\tau v^2/2\). Since \(\tau b/c<4/9\), the
derivative has opposite signs at \(\pm3b/2\) if \(b>0\), and curvature
stays positive between them. Thus its zero \(v_\circ\) satisfies
\(|v_\circ|\le3b/2\), also when \(b=0\). At that zero,
\(|v_\circ+g'(0)/c|\le9\tau b^2/(8c)\). This bounds the error using
\(\IFP^{-1}\nabla_t\). Replacing \(\IFP\) by \(\IF_t\) costs at most
\(\CONSTi_1\CONSTi_2/[\DFN_t(1-\delta_*)]\) in curvature units.
Use \(b\le\CONSTi_1/(1-\delta_*)\) and maximize over \(t\). The same localization gives
\(\max_t\DFN_t|\IMFt_t-\IMFs_t|\le3\CONSTi_1/[2(1-\delta_*)]\).
\end{proof}

\subsection{How the first error enters the later fits}\label{irmfproof:sec:recursive}

Take the Gaussian-smoothed absolute loss \(\rhohu=\rhohu_H\) of
Lemma~\ref{Lrhohu}. Write \(\KS_{tu}^{(k)}=\Kern_{h_k}(t-u)/N_k(t)\)
for the weights in \eqref{dfikkj3ujvyhwedrgg} and
\(\alpha_k=\max_{t,u}\KS_{tu}^{(k)}\).
The true component \(\IMFs^{(k)}\) uses the law of the actual residual
\(\Data^{(k)}\). For the proof we also need a deterministic reference:
start at \(\Signal_{\rhohum}^{(1)}=\Signal\), fit each row using
\(\rhohum(y)=\E\rhohu(y+\eps_1)\), then subtract that fit,
\[
 \IMFs_{\rhohum,t}^{(k)}
  =\arginf_s\sum_u\KS_{tu}^{(k)}\rhohum(\Signal_{\rhohum,u}^{(k)}-s),
 \qquad \Signal_{\rhohum}^{(k+1)}=\Signal_{\rhohum}^{(k)}-\IMFs_{\rhohum}^{(k)}.
\]
The two targets coincide at stage one. On the reference path, define
\(y_{ktu}=\Signal_{\rhohum,u}^{(k)}+\eps_u-\IMFs_{\rhohum,t}^{(k)}\),
\(g_{ktu}=\rhohu''(y_{ktu})\), and \(f_{kt}=\sum_u\KS_{tu}^{(k)}\E g_{ktu}\).
The score term \(b_k\) and the matrix \(V_k\) that transmits a change
in the input are
\[
 b_{kt}=\frac{\sum_u\KS_{tu}^{(k)}\rhohu'(y_{ktu})}{f_{kt}},
 \qquad V_k[t,u]=\frac{\KS_{tu}^{(k)}\E g_{ktu}}{f_{kt}} .
\]
In particular \(b_{1t}=\IF_t^{-1}\nabla_t\). At the next fit the leading
error is \(b_2-V_2b_1\). In general, accumulate the preceding leading
errors in \(\Lambda_k\) and set
\begin{equation}\label{irmfproof:eq:lambda}
 \Lambda_1=0,\qquad \mathcal L_k=b_k-V_k\Lambda_k,\qquad
 \Lambda_{k+1}=\Lambda_k+\mathcal L_k .
\end{equation}

\paragraph{Conditions.}
Assume \(f_{kt}\ge\lambda>0\) and
\(\alpha_j\le C_\alpha\alpha_k\hband^{-(k-j)}\) for \(j\le k\),
with \(\hband>1\), \(C_\alpha\ge1\). Write
\(B_{m:l}=(\Id-V_m)\cdots(\Id-V_l)\), with empty product \(\Id\).
We need stability of every such product and limited variation of curvature:
\begin{equation}\label{irmfproof:eq:dispersion}
 \max_{l\le m}\|B_{m:l}\|_\infty\le C_B,\qquad
 d=\max_{k,t}\frac{\sum_u\KS_{tu}^{(k)}\E|g_{ktu}-\E g_{ktu}|}{f_{kt}}
       <\frac{\hband-1}{C_B C_\alpha},\qquad C_B\ge1 .
\end{equation}
Here \(\|\cdot\|_\infty\) is the maximum absolute row sum for matrices and
the maximum absolute entry for vectors.

\begin{theorem}\label{irmfproof:thm:main}
Put \(u_k=\sqrt{z\alpha_k}\) and \(\alpha_{\max}=\max_k\alpha_k\).
There are constants \(C,\eta>0\) depending only on \(H,\lambda,\hband,
C_\alpha,C_B,d\) such that, if \(z\alpha_{\max}\le\eta\), then with
probability at least \(1-\ex^{-\xx}\), simultaneously for \(1\le k\le k_{\max}\),
\begin{equation}\label{irmfproof:eq:expansion}
 \|\mathcal L_k\|_\infty\le C u_k,\qquad
 \|\IMFt^{(k)}-\IMFs_{\rhohum}^{(k)}-\mathcal L_k\|_\infty\le C u_k^2 .
\end{equation}
\end{theorem}

\begin{proof}
Write \(e_k=\IMFt^{(k)}-\IMFs_{\rhohum}^{(k)}\) and
\(E_k=\sum_{j<k}e_j\). The actual input is exactly
\(\Data^{(k)}=\Signal_{\rhohum}^{(k)}+\eps-E_k\).
We use \(C\) for constants with the dependence stated above and put \(u_0=0\).

\emph{Expand the score once.}
The fitted location solves
\(\sum_u\KS_{tu}^{(k)}\rhohu'(y_{ktu}-E_k(u)-e_k(t))=0\).
Define \(Z_k[t,u]=\KS_{tu}^{(k)}(g_{ktu}-\E g_{ktu})/f_{kt}\).
Taylor expansion at \(y_{ktu}\), followed by division by \(f_{kt}\), gives
\begin{equation}\label{irmfproof:eq:recursion}
 e_k=b_k-V_kE_k+w_k,\qquad
 \|w_k+Z_kE_k\|_\infty
 \le \|Z_k\mathbf{1}\|_\infty\|e_k\|_\infty+
 \frac{L_3}{2\lambda}(\|E_k\|_\infty+\|e_k\|_\infty)^2 .
\end{equation}
Here \(Z_k\mathbf{1}\) is the row sum of the centered curvature. It is small
even when the sum of the absolute entries in that row is appreciable.

\emph{Control the terms that share observations.}
Since \(\sum_u(\KS_{tu}^{(k)})^2\le\alpha_k\), Bernstein's inequality
applied to the scores, curvatures and absolute centered curvatures gives,
simultaneously, \(\|b_k\|_\infty+\|Z_k\mathbf{1}\|_\infty\le C u_k\)
and \(\|Z_k\|_\infty\le d+C u_k\) when \(u_k\le1\).
For the input \(\Signal_{\rhohum}^{(k)}+\eps\), let \(s\) be the absolute
score at the reference fit. Here \(s\le C u_k\) and the empirical curvature
is at least \(3\lambda/4\) for small \(\eta\). Since it is \(L_3\)-Lipschitz,
it stays above \(\lambda/2\) within \(2s/\lambda\), which brackets the fit.
A local fit is monotone in each
observation and translation equivariant, hence 1-Lipschitz in the sup norm.
Consequently \(\|e_k\|_\infty\le C u_k+\|E_k\|_\infty\).

Iterating \eqref{irmfproof:eq:lambda} gives
\(\Lambda_k=\sum_{j<k}B_{k-1:j+1}b_j\). The geometric schedule then gives
\(\|\Lambda_k\|_\infty\le C_L u_{k-1}\).
The same sum grouped by observation gives
\(\Lambda_k(u)=\sum_v\gamma_{uv}(\eps_v)\), with each summand centered.
Indeed \(\gamma_{uv}\) is a sum of centered scores at \(v\), bounded in absolute value by
\(G_{uv}=2\lambda^{-1}\sum_{j<k}\sum_s|B_{k-1:j+1}[u,s]|\KS_{sv}^{(j)}\).
The product bound and the weight schedule imply
\(\max_{u,v}G_{uv}\le G_\infty\alpha_{k-1}\) and
\(\max_u\|G_{u,\cdot}\|_2\le C\sqrt{\alpha_{k-1}}\).
In \(Z_k\Lambda_k\), the terms \(u=v\) thus contribute at most
\(C\alpha_{k-1}\). The terms \(u\ne v\) are centered in each argument;
the sum of their squared envelopes is at most \(C\alpha_k\alpha_{k-1}\).
The quadratic-sum bound in the companion proof gives
\(\|Z_k\Lambda_k\|_\infty\le C u_k^2\).%
\footnote{For independent coordinates and separately centered kernels
\(|h_{uv}|\le H_{uv}\), the off-diagonal sum is at most
\(32\ex z(\sum H_{uv}^2)^{1/2}\) except with probability \(\ex^{-z}\)
for \(z\ge2\). A proof by symmetrization and explicit constants for the
component bounds are in Appendices A and B of the
\href{https://github.com/mkuziuk/imf/blob/ade89c8fd39ff9b3461ae68b33227ae08f9865a4/overleaf/irmf_all_components_bound.pdf}{companion proof}.}
The first three Bernstein bounds and this last bound cost at most
\(6n k_{\max}\ex^{-z}\le\ex^{-\xx}\) in total.

\emph{Close the induction.}
Set \(\Theta_k=E_k-\Lambda_k\). Choose
\(d<\delta<(\hband-1)/(C_B C_\alpha)\) and reduce \(\eta\) until
\(\|Z_k\|_\infty\le\delta\).
If \(\|\Theta_j\|_\infty\le C_\Theta u_{j-1}^2\) through stage \(k\), then
\(\|E_j\|_\infty\le(C_L+1)u_{j-1}\) provided \(C_\Theta\sqrt\eta\le1\).
Splitting \(Z_j E_j=Z_j\Lambda_j+Z_j\Theta_j\) in
\eqref{irmfproof:eq:recursion} now gives
\(\|w_j\|_\infty\le C_0 u_j^2+\delta C_\Theta u_{j-1}^2\), with \(C_0\)
independent of \(C_\Theta\). Since
\(\Theta_{k+1}=\sum_{j\le k}B_{k:j+1}w_j\), summing the two geometric
series gives
\[
 \|\Theta_{k+1}\|_\infty
 \le\left\{\frac{C_B C_\alpha C_0}{1-\hband^{-1}}
       +\frac{C_B C_\alpha\delta}{\hband-1}C_\Theta\right\}u_k^2 .
\]
The coefficient of \(C_\Theta\) is less than one. Choose \(C_\Theta\) large enough to
close this bound, then choose \(\eta\) as above. The induction starts at
\(E_1=\Lambda_1=0\). Finally \(e_k-\mathcal L_k=\Theta_{k+1}-\Theta_k\)
and \(u_{k-1}^2\le(C_\alpha/\hband)u_k^2\), which proves both bounds.
\end{proof}

\subsection{Returning to the true IMF}\label{irmfproof:sec:targets}

The first-stage target is already \(\IMFs_t\). At later stages the shift depends on how earlier fits changed the residuals.
In the expected score, contributions from distinct observations cancel.
Write \(z_0=\log(10n k_{\max})\), \(\xx_*=\eta/\alpha_{\max}-z_0>0\), and
assume \(\mu_4=\E|\eps_1|^4<\infty\). Let \(m_{\rhohu}\) solve
\(\E\rhohu'(\eps_1-m_{\rhohu})=0\) and put
\(M=2^{k_{\max}+1}(\|\Signal\|_\infty+(n\mu_4)^{1/4}+|m_{\rhohu}|)\).
With the induction constants above, set \(A=C_L+1\) and, for \(k\ge2\),
\begin{align*}
 U_k&=C_\Theta(z_0+1)\alpha_{k-1}+(M+C_B k_{\max}/\lambda)\ex^{-\xx_*/2},\\
 T_k&=A^2(z_0+1)\alpha_{k-1}+M^2\ex^{-\xx_*/2},\\
 \beta_k&=L_2(G_\infty\alpha_{k-1}+U_k)+L_3T_k/2,\qquad \beta_1=0 .
\end{align*}
\begin{corollary}\label{irmfproof:cor:population}
If \(L_3\sqrt{T_k}\le\lambda/4\) and \(\beta_k\le\lambda^2/(8L_3)\)
at each \(k\ge2\), the event of Theorem~\ref{irmfproof:thm:main} gives
\[
 \|\IMFt^{(k)}-\IMFs^{(k)}\|_\infty\le C(u_k+u_k^2)+2\beta_k/\lambda .
\]
If both exponential tail terms in \(U_k,T_k\) are at most \(\min_j\alpha_j\),
the shift is \(O(z\alpha_k)\). Thus for \(\alpha_k\asymp(nh_k)^{-1}\) the
simultaneous rate is \(\sqrt{z/(nh_k)}+z/(nh_k)\).
\end{corollary}

\begin{proof}
The induction gives \(\E\|\Theta_k\|_\infty\le U_k\) and
\(\E\|E_k\|_\infty^2\le T_k\): integrate its bounds over
\(0<\xx\le\xx_*\), using \(\int_0^\infty\ex^{-s}ds=1\).
Above the cutoff use Cauchy--Schwarz on an event of probability at most
\(\ex^{-\xx_*}\). Indeed an empirical fit lies between its extreme inputs
and a reference fit between those extremes shifted by \(m_{\rhohu}\);
iterating these brackets gives \(\|\|E_k\|_\infty\|_{L^4}\le M\).
Also \(\|\Lambda_k\|_\infty\le C_B k_{\max}/\lambda\) pathwise.

At \(\IMFs_{\rhohum,t}^{(k)}\) the true population score is
\(\sum_u\KS_{tu}^{(k)}\E\rhohu'(y_{ktu}-E_k(u))\).
Expand once about \(y_{ktu}\). The constant term sums to zero. In the linear term \(\E[g_{ktu}\Lambda_k(u)]\), independence and centering
cancel every \(\gamma_{uv}(\eps_v)\) except \(v=u\), leaving at most
\(L_2 G_\infty\alpha_{k-1}\). The remaining linear and quadratic terms
are bounded by \(L_2 U_k\) and \(L_3T_k/2\). Thus the score has magnitude
at most \(\beta_k\). The population loss has curvature at least
\(\lambda-L_3\sqrt{T_k}\) at the reference fit and at least \(\lambda/2\) within
\(2\beta_k/\lambda\). The score therefore has its zero within \(2\beta_k/\lambda\) of the reference fit.
Adding this displacement to \eqref{irmfproof:eq:expansion} proves the claim.
\end{proof}
 immediately after \input{IMF}.
\section{Finite-sample bounds for the IRMF components}\label{irmfproof:section}

At the first fit, the estimation error is the empirical score divided by
the population curvature, up to a smaller remainder. At the second fit,
part of that first error has already been subtracted from the data.
We show how to subtract its contribution from the second score and repeat
the argument at each scale. For a window containing \(m\) equally weighted
observations, the leading error at that scale is of order \(\sqrt{z/m}\)
and, under the conditions below, the remainder is of order \(z/m\), where \(z=\xx+\log(10n k_{\max})\)
at confidence \(1-\ex^{-\xx}\).

Throughout, the \(n\) grid points, the signal and the bandwidths are fixed,
\(\Data_u=\Signal_u+\eps_u\) with i.i.d.\ errors of finite first moment,
and the kernel weights are nonnegative with positive row sums. All maxima
over \(t\) are on this grid.

\subsection{The first component}\label{irmfproof:first}

In \eqref{miuygtd5r6wguwhd7r54we} let the contrast be convex, coercive and
\(C^3\), with \(|\rhohu'|\le L_1\), \(0\le\rhohu''\le L_2\), and
\(|\rhohu'''|\le L_3\). Write \(h=h_1\), \(Q_t^2=\sum_u\Kern_h(t-u)^2\),
\(\kappa_t=\max_u\Kern_h(t-u)\), \(\DFN_{\min}=\min_t\sqrt{\IF_t}>0\),
and \(z_1=\xx+\log(4n)\) for \(\xx>0\). The score, curvature and third
derivative in Lemma~\ref{LRMFloc} have respective bounds
\begin{align*}
 \CONSTi_1&=\max_t L_1(Q_t\sqrt{2z_1}+2\kappa_t z_1/3)/\DFN_t,\\
 \CONSTi_2&=\max_t L_2(Q_t\sqrt{2z_1}+\kappa_t z_1/3)/\DFN_t,
 &\CONSTi_3&=\max_t L_3N_h(t)/\DFN_t^2 .
\end{align*}

\begin{proposition}\label{irmfproof:first-bound}
If \(\delta_*=\CONSTi_2/\DFN_{\min}<1\) and
\(\CONSTi_1\CONSTi_3/[\DFN_{\min}(1-\delta_*)^2]<4/9\), then with
probability at least \(1-\ex^{-\xx}\),
\begin{equation}\label{irmfproof:first-expansion}
 \max_t\DFN_t|\IMFt_t-\IMFs_t-\IF_t^{-1}\nabla_t|
 \le \frac{1}{\DFN_{\min}}\left\{
 \frac{\CONSTi_1\CONSTi_2}{1-\delta_*}
 +\frac{9\CONSTi_1^2\CONSTi_3}{8(1-\delta_*)^3}\right\}.
\end{equation}
\end{proposition}

\begin{proof}
Fix \(t\), set \(\xs=\IMFs_t\), and use the empirical contrast \(\fp\)
in \eqref{fjiuy9iw3ojlggeefr}. The population score equation gives
\(\fp'(\xs)=-\nabla_t\). The centered score summands have total variance
at most \(L_1^2 Q_t^2\) and maximum size \(2L_1\kappa_t\); the corresponding
curvature bounds are \(L_2^2 Q_t^2\) and \(L_2\kappa_t\).
Bernstein's inequality gives \(|\nabla_t|\le\DFN_t\CONSTi_1\) and
\(|\IFP-\IF_t|\le\DFN_t\CONSTi_2\) for all \(t\), with failure probability
at most \(4n\ex^{-z_1}\), where \(\IFP=\fp''(\xs)\).
The third-derivative bound gives
\(\max(\sup|\fs'''|,\sup|\fp'''|)\le\DFN_t^2\CONSTi_3\), so
\(|\fs''(y)-\fs''(\xs)|\le\DFN_t^2\CONSTi_3|y-\xs|\).

In curvature units \(v=\DFN_t(y-\xs)\), put \(g(v)=\fp(y)\),
\(c=g''(0)\ge1-\delta_*\), \(b=|g'(0)|/c\), and
\(\tau=\CONSTi_3/\DFN_t\). Taylor's inequality gives
\(|g'(v)-g'(0)-cv|\le\tau v^2/2\). Since \(\tau b/c<4/9\), the
derivative has opposite signs at \(\pm3b/2\) if \(b>0\), and curvature
stays positive between them. Thus its zero \(v_\circ\) satisfies
\(|v_\circ|\le3b/2\), also when \(b=0\). At that zero,
\(|v_\circ+g'(0)/c|\le9\tau b^2/(8c)\). This bounds the error using
\(\IFP^{-1}\nabla_t\). Replacing \(\IFP\) by \(\IF_t\) costs at most
\(\CONSTi_1\CONSTi_2/[\DFN_t(1-\delta_*)]\) in curvature units.
Use \(b\le\CONSTi_1/(1-\delta_*)\) and maximize over \(t\). The same localization gives
\(\max_t\DFN_t|\IMFt_t-\IMFs_t|\le3\CONSTi_1/[2(1-\delta_*)]\).
\end{proof}

\subsection{How the first error enters the later fits}\label{irmfproof:sec:recursive}

Take the Gaussian-smoothed absolute loss \(\rhohu=\rhohu_H\) of
Lemma~\ref{Lrhohu}. Write \(\KS_{tu}^{(k)}=\Kern_{h_k}(t-u)/N_k(t)\)
for the weights in \eqref{dfikkj3ujvyhwedrgg} and
\(\alpha_k=\max_{t,u}\KS_{tu}^{(k)}\).
The true component \(\IMFs^{(k)}\) uses the law of the actual residual
\(\Data^{(k)}\). For the proof we also need a deterministic reference:
start at \(\Signal_{\rhohum}^{(1)}=\Signal\), fit each row using
\(\rhohum(y)=\E\rhohu(y+\eps_1)\), then subtract that fit,
\[
 \IMFs_{\rhohum,t}^{(k)}
  =\arginf_s\sum_u\KS_{tu}^{(k)}\rhohum(\Signal_{\rhohum,u}^{(k)}-s),
 \qquad \Signal_{\rhohum}^{(k+1)}=\Signal_{\rhohum}^{(k)}-\IMFs_{\rhohum}^{(k)}.
\]
The two targets coincide at stage one. On the reference path, define
\(y_{ktu}=\Signal_{\rhohum,u}^{(k)}+\eps_u-\IMFs_{\rhohum,t}^{(k)}\),
\(g_{ktu}=\rhohu''(y_{ktu})\), and \(f_{kt}=\sum_u\KS_{tu}^{(k)}\E g_{ktu}\).
The score term \(b_k\) and the matrix \(V_k\) that transmits a change
in the input are
\[
 b_{kt}=\frac{\sum_u\KS_{tu}^{(k)}\rhohu'(y_{ktu})}{f_{kt}},
 \qquad V_k[t,u]=\frac{\KS_{tu}^{(k)}\E g_{ktu}}{f_{kt}} .
\]
In particular \(b_{1t}=\IF_t^{-1}\nabla_t\). At the next fit the leading
error is \(b_2-V_2b_1\). In general, accumulate the preceding leading
errors in \(\Lambda_k\) and set
\begin{equation}\label{irmfproof:eq:lambda}
 \Lambda_1=0,\qquad \mathcal L_k=b_k-V_k\Lambda_k,\qquad
 \Lambda_{k+1}=\Lambda_k+\mathcal L_k .
\end{equation}

\paragraph{Conditions.}
Assume \(f_{kt}\ge\lambda>0\) and
\(\alpha_j\le C_\alpha\alpha_k\hband^{-(k-j)}\) for \(j\le k\),
with \(\hband>1\), \(C_\alpha\ge1\). Write
\(B_{m:l}=(\Id-V_m)\cdots(\Id-V_l)\), with empty product \(\Id\).
We need stability of every such product and limited variation of curvature:
\begin{equation}\label{irmfproof:eq:dispersion}
 \max_{l\le m}\|B_{m:l}\|_\infty\le C_B,\qquad
 d=\max_{k,t}\frac{\sum_u\KS_{tu}^{(k)}\E|g_{ktu}-\E g_{ktu}|}{f_{kt}}
       <\frac{\hband-1}{C_B C_\alpha},\qquad C_B\ge1 .
\end{equation}
Here \(\|\cdot\|_\infty\) is the maximum absolute row sum for matrices and
the maximum absolute entry for vectors.

\begin{theorem}\label{irmfproof:thm:main}
Put \(u_k=\sqrt{z\alpha_k}\) and \(\alpha_{\max}=\max_k\alpha_k\).
There are constants \(C,\eta>0\) depending only on \(H,\lambda,\hband,
C_\alpha,C_B,d\) such that, if \(z\alpha_{\max}\le\eta\), then with
probability at least \(1-\ex^{-\xx}\), simultaneously for \(1\le k\le k_{\max}\),
\begin{equation}\label{irmfproof:eq:expansion}
 \|\mathcal L_k\|_\infty\le C u_k,\qquad
 \|\IMFt^{(k)}-\IMFs_{\rhohum}^{(k)}-\mathcal L_k\|_\infty\le C u_k^2 .
\end{equation}
\end{theorem}

\begin{proof}
Write \(e_k=\IMFt^{(k)}-\IMFs_{\rhohum}^{(k)}\) and
\(E_k=\sum_{j<k}e_j\). The actual input is exactly
\(\Data^{(k)}=\Signal_{\rhohum}^{(k)}+\eps-E_k\).
We use \(C\) for constants with the dependence stated above and put \(u_0=0\).

\emph{Expand the score once.}
The fitted location solves
\(\sum_u\KS_{tu}^{(k)}\rhohu'(y_{ktu}-E_k(u)-e_k(t))=0\).
Define \(Z_k[t,u]=\KS_{tu}^{(k)}(g_{ktu}-\E g_{ktu})/f_{kt}\).
Taylor expansion at \(y_{ktu}\), followed by division by \(f_{kt}\), gives
\begin{equation}\label{irmfproof:eq:recursion}
 e_k=b_k-V_kE_k+w_k,\qquad
 \|w_k+Z_kE_k\|_\infty
 \le \|Z_k\mathbf{1}\|_\infty\|e_k\|_\infty+
 \frac{L_3}{2\lambda}(\|E_k\|_\infty+\|e_k\|_\infty)^2 .
\end{equation}
Here \(Z_k\mathbf{1}\) is the row sum of the centered curvature. It is small
even when the sum of the absolute entries in that row is appreciable.

\emph{Control the terms that share observations.}
Since \(\sum_u(\KS_{tu}^{(k)})^2\le\alpha_k\), Bernstein's inequality
applied to the scores, curvatures and absolute centered curvatures gives,
simultaneously, \(\|b_k\|_\infty+\|Z_k\mathbf{1}\|_\infty\le C u_k\)
and \(\|Z_k\|_\infty\le d+C u_k\) when \(u_k\le1\).
For the input \(\Signal_{\rhohum}^{(k)}+\eps\), let \(s\) be the absolute
score at the reference fit. Here \(s\le C u_k\) and the empirical curvature
is at least \(3\lambda/4\) for small \(\eta\). Since it is \(L_3\)-Lipschitz,
it stays above \(\lambda/2\) within \(2s/\lambda\), which brackets the fit.
A local fit is monotone in each
observation and translation equivariant, hence 1-Lipschitz in the sup norm.
Consequently \(\|e_k\|_\infty\le C u_k+\|E_k\|_\infty\).

Iterating \eqref{irmfproof:eq:lambda} gives
\(\Lambda_k=\sum_{j<k}B_{k-1:j+1}b_j\). The geometric schedule then gives
\(\|\Lambda_k\|_\infty\le C_L u_{k-1}\).
The same sum grouped by observation gives
\(\Lambda_k(u)=\sum_v\gamma_{uv}(\eps_v)\), with each summand centered.
Indeed \(\gamma_{uv}\) is a sum of centered scores at \(v\), bounded in absolute value by
\(G_{uv}=2\lambda^{-1}\sum_{j<k}\sum_s|B_{k-1:j+1}[u,s]|\KS_{sv}^{(j)}\).
The product bound and the weight schedule imply
\(\max_{u,v}G_{uv}\le G_\infty\alpha_{k-1}\) and
\(\max_u\|G_{u,\cdot}\|_2\le C\sqrt{\alpha_{k-1}}\).
In \(Z_k\Lambda_k\), the terms \(u=v\) thus contribute at most
\(C\alpha_{k-1}\). The terms \(u\ne v\) are centered in each argument;
the sum of their squared envelopes is at most \(C\alpha_k\alpha_{k-1}\).
The quadratic-sum bound in the companion proof gives
\(\|Z_k\Lambda_k\|_\infty\le C u_k^2\).%
\footnote{For independent coordinates and separately centered kernels
\(|h_{uv}|\le H_{uv}\), the off-diagonal sum is at most
\(32\ex z(\sum H_{uv}^2)^{1/2}\) except with probability \(\ex^{-z}\)
for \(z\ge2\). A proof by symmetrization and explicit constants for the
component bounds are in Appendices A and B of the
\href{https://github.com/mkuziuk/imf/blob/ade89c8fd39ff9b3461ae68b33227ae08f9865a4/overleaf/irmf_all_components_bound.pdf}{companion proof}.}
The first three Bernstein bounds and this last bound cost at most
\(6n k_{\max}\ex^{-z}\le\ex^{-\xx}\) in total.

\emph{Close the induction.}
Set \(\Theta_k=E_k-\Lambda_k\). Choose
\(d<\delta<(\hband-1)/(C_B C_\alpha)\) and reduce \(\eta\) until
\(\|Z_k\|_\infty\le\delta\).
If \(\|\Theta_j\|_\infty\le C_\Theta u_{j-1}^2\) through stage \(k\), then
\(\|E_j\|_\infty\le(C_L+1)u_{j-1}\) provided \(C_\Theta\sqrt\eta\le1\).
Splitting \(Z_j E_j=Z_j\Lambda_j+Z_j\Theta_j\) in
\eqref{irmfproof:eq:recursion} now gives
\(\|w_j\|_\infty\le C_0 u_j^2+\delta C_\Theta u_{j-1}^2\), with \(C_0\)
independent of \(C_\Theta\). Since
\(\Theta_{k+1}=\sum_{j\le k}B_{k:j+1}w_j\), summing the two geometric
series gives
\[
 \|\Theta_{k+1}\|_\infty
 \le\left\{\frac{C_B C_\alpha C_0}{1-\hband^{-1}}
       +\frac{C_B C_\alpha\delta}{\hband-1}C_\Theta\right\}u_k^2 .
\]
The coefficient of \(C_\Theta\) is less than one. Choose \(C_\Theta\) large enough to
close this bound, then choose \(\eta\) as above. The induction starts at
\(E_1=\Lambda_1=0\). Finally \(e_k-\mathcal L_k=\Theta_{k+1}-\Theta_k\)
and \(u_{k-1}^2\le(C_\alpha/\hband)u_k^2\), which proves both bounds.
\end{proof}

\subsection{Returning to the true IMF}\label{irmfproof:sec:targets}

The first-stage target is already \(\IMFs_t\). At later stages the shift depends on how earlier fits changed the residuals.
In the expected score, contributions from distinct observations cancel.
Write \(z_0=\log(10n k_{\max})\), \(\xx_*=\eta/\alpha_{\max}-z_0>0\), and
assume \(\mu_4=\E|\eps_1|^4<\infty\). Let \(m_{\rhohu}\) solve
\(\E\rhohu'(\eps_1-m_{\rhohu})=0\) and put
\(M=2^{k_{\max}+1}(\|\Signal\|_\infty+(n\mu_4)^{1/4}+|m_{\rhohu}|)\).
With the induction constants above, set \(A=C_L+1\) and, for \(k\ge2\),
\begin{align*}
 U_k&=C_\Theta(z_0+1)\alpha_{k-1}+(M+C_B k_{\max}/\lambda)\ex^{-\xx_*/2},\\
 T_k&=A^2(z_0+1)\alpha_{k-1}+M^2\ex^{-\xx_*/2},\\
 \beta_k&=L_2(G_\infty\alpha_{k-1}+U_k)+L_3T_k/2,\qquad \beta_1=0 .
\end{align*}
\begin{corollary}\label{irmfproof:cor:population}
If \(L_3\sqrt{T_k}\le\lambda/4\) and \(\beta_k\le\lambda^2/(8L_3)\)
at each \(k\ge2\), the event of Theorem~\ref{irmfproof:thm:main} gives
\[
 \|\IMFt^{(k)}-\IMFs^{(k)}\|_\infty\le C(u_k+u_k^2)+2\beta_k/\lambda .
\]
If both exponential tail terms in \(U_k,T_k\) are at most \(\min_j\alpha_j\),
the shift is \(O(z\alpha_k)\). Thus for \(\alpha_k\asymp(nh_k)^{-1}\) the
simultaneous rate is \(\sqrt{z/(nh_k)}+z/(nh_k)\).
\end{corollary}

\begin{proof}
The induction gives \(\E\|\Theta_k\|_\infty\le U_k\) and
\(\E\|E_k\|_\infty^2\le T_k\): integrate its bounds over
\(0<\xx\le\xx_*\), using \(\int_0^\infty\ex^{-s}ds=1\).
Above the cutoff use Cauchy--Schwarz on an event of probability at most
\(\ex^{-\xx_*}\). Indeed an empirical fit lies between its extreme inputs
and a reference fit between those extremes shifted by \(m_{\rhohu}\);
iterating these brackets gives \(\|\|E_k\|_\infty\|_{L^4}\le M\).
Also \(\|\Lambda_k\|_\infty\le C_B k_{\max}/\lambda\) pathwise.

At \(\IMFs_{\rhohum,t}^{(k)}\) the true population score is
\(\sum_u\KS_{tu}^{(k)}\E\rhohu'(y_{ktu}-E_k(u))\).
Expand once about \(y_{ktu}\). The constant term sums to zero. In the linear term \(\E[g_{ktu}\Lambda_k(u)]\), independence and centering
cancel every \(\gamma_{uv}(\eps_v)\) except \(v=u\), leaving at most
\(L_2 G_\infty\alpha_{k-1}\). The remaining linear and quadratic terms
are bounded by \(L_2 U_k\) and \(L_3T_k/2\). Thus the score has magnitude
at most \(\beta_k\). The population loss has curvature at least
\(\lambda-L_3\sqrt{T_k}\) at the reference fit and at least \(\lambda/2\) within
\(2\beta_k/\lambda\). The score therefore has its zero within \(2\beta_k/\lambda\) of the reference fit.
Adding this displacement to \eqref{irmfproof:eq:expansion} proves the claim.
\end{proof}
 immediately after \input{IMF}.
\section{Finite-sample bounds for the IRMF components}\label{irmfproof:section}

At the first fit, the estimation error is the empirical score divided by
the population curvature, up to a smaller remainder. At the second fit,
part of that first error has already been subtracted from the data.
We show how to subtract its contribution from the second score and repeat
the argument at each scale. For a window containing \(m\) equally weighted
observations, the leading error at that scale is of order \(\sqrt{z/m}\)
and, under the conditions below, the remainder is of order \(z/m\), where \(z=\xx+\log(10n k_{\max})\)
at confidence \(1-\ex^{-\xx}\).

Throughout, the \(n\) grid points, the signal and the bandwidths are fixed,
\(\Data_u=\Signal_u+\eps_u\) with i.i.d.\ errors of finite first moment,
and the kernel weights are nonnegative with positive row sums. All maxima
over \(t\) are on this grid.

\subsection{The first component}\label{irmfproof:first}

In \eqref{miuygtd5r6wguwhd7r54we} let the contrast be convex, coercive and
\(C^3\), with \(|\rhohu'|\le L_1\), \(0\le\rhohu''\le L_2\), and
\(|\rhohu'''|\le L_3\). Write \(h=h_1\), \(Q_t^2=\sum_u\Kern_h(t-u)^2\),
\(\kappa_t=\max_u\Kern_h(t-u)\), \(\DFN_{\min}=\min_t\sqrt{\IF_t}>0\),
and \(z_1=\xx+\log(4n)\) for \(\xx>0\). The score, curvature and third
derivative in Lemma~\ref{LRMFloc} have respective bounds
\begin{align*}
 \CONSTi_1&=\max_t L_1(Q_t\sqrt{2z_1}+2\kappa_t z_1/3)/\DFN_t,\\
 \CONSTi_2&=\max_t L_2(Q_t\sqrt{2z_1}+\kappa_t z_1/3)/\DFN_t,
 &\CONSTi_3&=\max_t L_3N_h(t)/\DFN_t^2 .
\end{align*}

\begin{proposition}\label{irmfproof:first-bound}
If \(\delta_*=\CONSTi_2/\DFN_{\min}<1\) and
\(\CONSTi_1\CONSTi_3/[\DFN_{\min}(1-\delta_*)^2]<4/9\), then with
probability at least \(1-\ex^{-\xx}\),
\begin{equation}\label{irmfproof:first-expansion}
 \max_t\DFN_t|\IMFt_t-\IMFs_t-\IF_t^{-1}\nabla_t|
 \le \frac{1}{\DFN_{\min}}\left\{
 \frac{\CONSTi_1\CONSTi_2}{1-\delta_*}
 +\frac{9\CONSTi_1^2\CONSTi_3}{8(1-\delta_*)^3}\right\}.
\end{equation}
\end{proposition}

\begin{proof}
Fix \(t\), set \(\xs=\IMFs_t\), and use the empirical contrast \(\fp\)
in \eqref{fjiuy9iw3ojlggeefr}. The population score equation gives
\(\fp'(\xs)=-\nabla_t\). The centered score summands have total variance
at most \(L_1^2 Q_t^2\) and maximum size \(2L_1\kappa_t\); the corresponding
curvature bounds are \(L_2^2 Q_t^2\) and \(L_2\kappa_t\).
Bernstein's inequality gives \(|\nabla_t|\le\DFN_t\CONSTi_1\) and
\(|\IFP-\IF_t|\le\DFN_t\CONSTi_2\) for all \(t\), with failure probability
at most \(4n\ex^{-z_1}\), where \(\IFP=\fp''(\xs)\).
The third-derivative bound gives
\(\max(\sup|\fs'''|,\sup|\fp'''|)\le\DFN_t^2\CONSTi_3\), so
\(|\fs''(y)-\fs''(\xs)|\le\DFN_t^2\CONSTi_3|y-\xs|\).

In curvature units \(v=\DFN_t(y-\xs)\), put \(g(v)=\fp(y)\),
\(c=g''(0)\ge1-\delta_*\), \(b=|g'(0)|/c\), and
\(\tau=\CONSTi_3/\DFN_t\). Taylor's inequality gives
\(|g'(v)-g'(0)-cv|\le\tau v^2/2\). Since \(\tau b/c<4/9\), the
derivative has opposite signs at \(\pm3b/2\) if \(b>0\), and curvature
stays positive between them. Thus its zero \(v_\circ\) satisfies
\(|v_\circ|\le3b/2\), also when \(b=0\). At that zero,
\(|v_\circ+g'(0)/c|\le9\tau b^2/(8c)\). This bounds the error using
\(\IFP^{-1}\nabla_t\). Replacing \(\IFP\) by \(\IF_t\) costs at most
\(\CONSTi_1\CONSTi_2/[\DFN_t(1-\delta_*)]\) in curvature units.
Use \(b\le\CONSTi_1/(1-\delta_*)\) and maximize over \(t\). The same localization gives
\(\max_t\DFN_t|\IMFt_t-\IMFs_t|\le3\CONSTi_1/[2(1-\delta_*)]\).
\end{proof}

\subsection{How the first error enters the later fits}\label{irmfproof:sec:recursive}

Take the Gaussian-smoothed absolute loss \(\rhohu=\rhohu_H\) of
Lemma~\ref{Lrhohu}. Write \(\KS_{tu}^{(k)}=\Kern_{h_k}(t-u)/N_k(t)\)
for the weights in \eqref{dfikkj3ujvyhwedrgg} and
\(\alpha_k=\max_{t,u}\KS_{tu}^{(k)}\).
The true component \(\IMFs^{(k)}\) uses the law of the actual residual
\(\Data^{(k)}\). For the proof we also need a deterministic reference:
start at \(\Signal_{\rhohum}^{(1)}=\Signal\), fit each row using
\(\rhohum(y)=\E\rhohu(y+\eps_1)\), then subtract that fit,
\[
 \IMFs_{\rhohum,t}^{(k)}
  =\arginf_s\sum_u\KS_{tu}^{(k)}\rhohum(\Signal_{\rhohum,u}^{(k)}-s),
 \qquad \Signal_{\rhohum}^{(k+1)}=\Signal_{\rhohum}^{(k)}-\IMFs_{\rhohum}^{(k)}.
\]
The two targets coincide at stage one. On the reference path, define
\(y_{ktu}=\Signal_{\rhohum,u}^{(k)}+\eps_u-\IMFs_{\rhohum,t}^{(k)}\),
\(g_{ktu}=\rhohu''(y_{ktu})\), and \(f_{kt}=\sum_u\KS_{tu}^{(k)}\E g_{ktu}\).
The score term \(b_k\) and the matrix \(V_k\) that transmits a change
in the input are
\[
 b_{kt}=\frac{\sum_u\KS_{tu}^{(k)}\rhohu'(y_{ktu})}{f_{kt}},
 \qquad V_k[t,u]=\frac{\KS_{tu}^{(k)}\E g_{ktu}}{f_{kt}} .
\]
In particular \(b_{1t}=\IF_t^{-1}\nabla_t\). At the next fit the leading
error is \(b_2-V_2b_1\). In general, accumulate the preceding leading
errors in \(\Lambda_k\) and set
\begin{equation}\label{irmfproof:eq:lambda}
 \Lambda_1=0,\qquad \mathcal L_k=b_k-V_k\Lambda_k,\qquad
 \Lambda_{k+1}=\Lambda_k+\mathcal L_k .
\end{equation}

\paragraph{Conditions.}
Assume \(f_{kt}\ge\lambda>0\) and
\(\alpha_j\le C_\alpha\alpha_k\hband^{-(k-j)}\) for \(j\le k\),
with \(\hband>1\), \(C_\alpha\ge1\). Write
\(B_{m:l}=(\Id-V_m)\cdots(\Id-V_l)\), with empty product \(\Id\).
We need stability of every such product and limited variation of curvature:
\begin{equation}\label{irmfproof:eq:dispersion}
 \max_{l\le m}\|B_{m:l}\|_\infty\le C_B,\qquad
 d=\max_{k,t}\frac{\sum_u\KS_{tu}^{(k)}\E|g_{ktu}-\E g_{ktu}|}{f_{kt}}
       <\frac{\hband-1}{C_B C_\alpha},\qquad C_B\ge1 .
\end{equation}
Here \(\|\cdot\|_\infty\) is the maximum absolute row sum for matrices and
the maximum absolute entry for vectors.

\begin{theorem}\label{irmfproof:thm:main}
Put \(u_k=\sqrt{z\alpha_k}\) and \(\alpha_{\max}=\max_k\alpha_k\).
There are constants \(C,\eta>0\) depending only on \(H,\lambda,\hband,
C_\alpha,C_B,d\) such that, if \(z\alpha_{\max}\le\eta\), then with
probability at least \(1-\ex^{-\xx}\), simultaneously for \(1\le k\le k_{\max}\),
\begin{equation}\label{irmfproof:eq:expansion}
 \|\mathcal L_k\|_\infty\le C u_k,\qquad
 \|\IMFt^{(k)}-\IMFs_{\rhohum}^{(k)}-\mathcal L_k\|_\infty\le C u_k^2 .
\end{equation}
\end{theorem}

\begin{proof}
Write \(e_k=\IMFt^{(k)}-\IMFs_{\rhohum}^{(k)}\) and
\(E_k=\sum_{j<k}e_j\). The actual input is exactly
\(\Data^{(k)}=\Signal_{\rhohum}^{(k)}+\eps-E_k\).
We use \(C\) for constants with the dependence stated above and put \(u_0=0\).

\emph{Expand the score once.}
The fitted location solves
\(\sum_u\KS_{tu}^{(k)}\rhohu'(y_{ktu}-E_k(u)-e_k(t))=0\).
Define \(Z_k[t,u]=\KS_{tu}^{(k)}(g_{ktu}-\E g_{ktu})/f_{kt}\).
Taylor expansion at \(y_{ktu}\), followed by division by \(f_{kt}\), gives
\begin{equation}\label{irmfproof:eq:recursion}
 e_k=b_k-V_kE_k+w_k,\qquad
 \|w_k+Z_kE_k\|_\infty
 \le \|Z_k\mathbf{1}\|_\infty\|e_k\|_\infty+
 \frac{L_3}{2\lambda}(\|E_k\|_\infty+\|e_k\|_\infty)^2 .
\end{equation}
Here \(Z_k\mathbf{1}\) is the row sum of the centered curvature. It is small
even when the sum of the absolute entries in that row is appreciable.

\emph{Control the terms that share observations.}
Since \(\sum_u(\KS_{tu}^{(k)})^2\le\alpha_k\), Bernstein's inequality
applied to the scores, curvatures and absolute centered curvatures gives,
simultaneously, \(\|b_k\|_\infty+\|Z_k\mathbf{1}\|_\infty\le C u_k\)
and \(\|Z_k\|_\infty\le d+C u_k\) when \(u_k\le1\).
For the input \(\Signal_{\rhohum}^{(k)}+\eps\), let \(s\) be the absolute
score at the reference fit. Here \(s\le C u_k\) and the empirical curvature
is at least \(3\lambda/4\) for small \(\eta\). Since it is \(L_3\)-Lipschitz,
it stays above \(\lambda/2\) within \(2s/\lambda\), which brackets the fit.
A local fit is monotone in each
observation and translation equivariant, hence 1-Lipschitz in the sup norm.
Consequently \(\|e_k\|_\infty\le C u_k+\|E_k\|_\infty\).

Iterating \eqref{irmfproof:eq:lambda} gives
\(\Lambda_k=\sum_{j<k}B_{k-1:j+1}b_j\). The geometric schedule then gives
\(\|\Lambda_k\|_\infty\le C_L u_{k-1}\).
The same sum grouped by observation gives
\(\Lambda_k(u)=\sum_v\gamma_{uv}(\eps_v)\), with each summand centered.
Indeed \(\gamma_{uv}\) is a sum of centered scores at \(v\), bounded in absolute value by
\(G_{uv}=2\lambda^{-1}\sum_{j<k}\sum_s|B_{k-1:j+1}[u,s]|\KS_{sv}^{(j)}\).
The product bound and the weight schedule imply
\(\max_{u,v}G_{uv}\le G_\infty\alpha_{k-1}\) and
\(\max_u\|G_{u,\cdot}\|_2\le C\sqrt{\alpha_{k-1}}\).
In \(Z_k\Lambda_k\), the terms \(u=v\) thus contribute at most
\(C\alpha_{k-1}\). The terms \(u\ne v\) are centered in each argument;
the sum of their squared envelopes is at most \(C\alpha_k\alpha_{k-1}\).
The quadratic-sum bound in the companion proof gives
\(\|Z_k\Lambda_k\|_\infty\le C u_k^2\).%
\footnote{For independent coordinates and separately centered kernels
\(|h_{uv}|\le H_{uv}\), the off-diagonal sum is at most
\(32\ex z(\sum H_{uv}^2)^{1/2}\) except with probability \(\ex^{-z}\)
for \(z\ge2\). A proof by symmetrization and explicit constants for the
component bounds are in Appendices A and B of the
\href{https://github.com/mkuziuk/imf/blob/ade89c8fd39ff9b3461ae68b33227ae08f9865a4/overleaf/irmf_all_components_bound.pdf}{companion proof}.}
The first three Bernstein bounds and this last bound cost at most
\(6n k_{\max}\ex^{-z}\le\ex^{-\xx}\) in total.

\emph{Close the induction.}
Set \(\Theta_k=E_k-\Lambda_k\). Choose
\(d<\delta<(\hband-1)/(C_B C_\alpha)\) and reduce \(\eta\) until
\(\|Z_k\|_\infty\le\delta\).
If \(\|\Theta_j\|_\infty\le C_\Theta u_{j-1}^2\) through stage \(k\), then
\(\|E_j\|_\infty\le(C_L+1)u_{j-1}\) provided \(C_\Theta\sqrt\eta\le1\).
Splitting \(Z_j E_j=Z_j\Lambda_j+Z_j\Theta_j\) in
\eqref{irmfproof:eq:recursion} now gives
\(\|w_j\|_\infty\le C_0 u_j^2+\delta C_\Theta u_{j-1}^2\), with \(C_0\)
independent of \(C_\Theta\). Since
\(\Theta_{k+1}=\sum_{j\le k}B_{k:j+1}w_j\), summing the two geometric
series gives
\[
 \|\Theta_{k+1}\|_\infty
 \le\left\{\frac{C_B C_\alpha C_0}{1-\hband^{-1}}
       +\frac{C_B C_\alpha\delta}{\hband-1}C_\Theta\right\}u_k^2 .
\]
The coefficient of \(C_\Theta\) is less than one. Choose \(C_\Theta\) large enough to
close this bound, then choose \(\eta\) as above. The induction starts at
\(E_1=\Lambda_1=0\). Finally \(e_k-\mathcal L_k=\Theta_{k+1}-\Theta_k\)
and \(u_{k-1}^2\le(C_\alpha/\hband)u_k^2\), which proves both bounds.
\end{proof}

\subsection{Returning to the true IMF}\label{irmfproof:sec:targets}

The first-stage target is already \(\IMFs_t\). At later stages the shift depends on how earlier fits changed the residuals.
In the expected score, contributions from distinct observations cancel.
Write \(z_0=\log(10n k_{\max})\), \(\xx_*=\eta/\alpha_{\max}-z_0>0\), and
assume \(\mu_4=\E|\eps_1|^4<\infty\). Let \(m_{\rhohu}\) solve
\(\E\rhohu'(\eps_1-m_{\rhohu})=0\) and put
\(M=2^{k_{\max}+1}(\|\Signal\|_\infty+(n\mu_4)^{1/4}+|m_{\rhohu}|)\).
With the induction constants above, set \(A=C_L+1\) and, for \(k\ge2\),
\begin{align*}
 U_k&=C_\Theta(z_0+1)\alpha_{k-1}+(M+C_B k_{\max}/\lambda)\ex^{-\xx_*/2},\\
 T_k&=A^2(z_0+1)\alpha_{k-1}+M^2\ex^{-\xx_*/2},\\
 \beta_k&=L_2(G_\infty\alpha_{k-1}+U_k)+L_3T_k/2,\qquad \beta_1=0 .
\end{align*}
\begin{corollary}\label{irmfproof:cor:population}
If \(L_3\sqrt{T_k}\le\lambda/4\) and \(\beta_k\le\lambda^2/(8L_3)\)
at each \(k\ge2\), the event of Theorem~\ref{irmfproof:thm:main} gives
\[
 \|\IMFt^{(k)}-\IMFs^{(k)}\|_\infty\le C(u_k+u_k^2)+2\beta_k/\lambda .
\]
If both exponential tail terms in \(U_k,T_k\) are at most \(\min_j\alpha_j\),
the shift is \(O(z\alpha_k)\). Thus for \(\alpha_k\asymp(nh_k)^{-1}\) the
simultaneous rate is \(\sqrt{z/(nh_k)}+z/(nh_k)\).
\end{corollary}

\begin{proof}
The induction gives \(\E\|\Theta_k\|_\infty\le U_k\) and
\(\E\|E_k\|_\infty^2\le T_k\): integrate its bounds over
\(0<\xx\le\xx_*\), using \(\int_0^\infty\ex^{-s}ds=1\).
Above the cutoff use Cauchy--Schwarz on an event of probability at most
\(\ex^{-\xx_*}\). Indeed an empirical fit lies between its extreme inputs
and a reference fit between those extremes shifted by \(m_{\rhohu}\);
iterating these brackets gives \(\|\|E_k\|_\infty\|_{L^4}\le M\).
Also \(\|\Lambda_k\|_\infty\le C_B k_{\max}/\lambda\) pathwise.

At \(\IMFs_{\rhohum,t}^{(k)}\) the true population score is
\(\sum_u\KS_{tu}^{(k)}\E\rhohu'(y_{ktu}-E_k(u))\).
Expand once about \(y_{ktu}\). The constant term sums to zero. In the linear term \(\E[g_{ktu}\Lambda_k(u)]\), independence and centering
cancel every \(\gamma_{uv}(\eps_v)\) except \(v=u\), leaving at most
\(L_2 G_\infty\alpha_{k-1}\). The remaining linear and quadratic terms
are bounded by \(L_2 U_k\) and \(L_3T_k/2\). Thus the score has magnitude
at most \(\beta_k\). The population loss has curvature at least
\(\lambda-L_3\sqrt{T_k}\) at the reference fit and at least \(\lambda/2\) within
\(2\beta_k/\lambda\). The score therefore has its zero within \(2\beta_k/\lambda\) of the reference fit.
Adding this displacement to \eqref{irmfproof:eq:expansion} proves the claim.
\end{proof}
